% =====================================================================
% SECTION 4 — THEORETICAL PROPERTIES (strengthened, standalone)
% Reuses existing label sec:theory and proposition prop:q-validity.
% New numbered propositions: consistency/MDE, local power, orthogonality,
% partial identification / measure invariance.
% All numerics verbatim from paper.tex; all \cite keys exist in thebibliography.
% =====================================================================

\section{Theoretical properties}
\label{sec:theory}

This section collects the formal properties of the test. We establish four
things: that the Monte Carlo test is finite-sample exact under
$\mathbb{Q}_\theta$-exchangeability; that it is consistent, with a minimum
detectable edge of order $N^{-1/2}$ in the trade count; that it attains the
local power of the one-sided Gaussian test against placement alternatives while
carrying \emph{zero} noncentrality against the pure exposure direction; and that,
because the conditioning measure is a modeling choice, the timing estimand is
only partially identified---the data return a set of verdicts over the
admissible measure class, and a valid test exists for the boundary of
measure-invariant skill. Throughout, $T(S)$ is the cumulative-return statistic
induced by a schedule $S$, $S_0$ is the realized schedule, and $\mathcal C=
\sigma(\mathcal S,\{P_t\})$ is the conditioning $\sigma$-field of
Section~\ref{sec:conditioning-estimand}.

\subsection{Finite-sample exactness}

The validity of the procedure does not rest on any model for returns; it rests
on the exchangeability of schedules under the resampling law. We restate the
result so that the section is self-contained.

\begin{proposition}[Finite-sample validity and large-$M$ consistency of the
Monte Carlo $p$-value]
\label{prop:q-validity}
Fix a measure $\mathbb{Q}_\theta$ in the family of Definition~\ref{def:h0}, and
let $S_1,\dots,S_M$ be i.i.d.\ draws from $\mathbb{Q}_\theta$ given $\mathcal C$.
Define
\[
\hat p_M
=\frac{1+\sum_{m=1}^{M}\mathbf 1\{T(S_m)\ge T(S_0)\}}{M+1}.
\]
\emph{(i) Validity.} Under $H_0^\theta$, that $S_0$ is itself a draw from
$\mathbb{Q}_\theta$ given $\mathcal C$, the statistic $\hat p_M$ is super-uniform:
for every $\alpha$ on the grid $\{1/(M+1),\dots,1\}$,
$\Pr(\hat p_M\le\alpha\mid\mathcal C)\le\alpha$, with equality absent ties and
strict conservatism when ties are counted into the right tail.
\emph{(ii) Consistency in $M$.} For \emph{any} fixed $S_0$, whether or not
$H_0^\theta$ holds, $\hat p_M\to q_\theta:=\mathbb{Q}_\theta\big(T(S)\ge
T(S_0)\mid\mathcal C\big)$ almost surely as $M\to\infty$.
\end{proposition}

\begin{proof}
(i) Under $H_0^\theta$ the variables $S_0,S_1,\dots,S_M$ are exchangeable given
$\mathcal C$, so the rank of $T(S_0)$ among $\{T(S_m)\}_{m=0}^{M}$ is uniform on
$\{1,\dots,M+1\}$ absent ties. Dividing the upper-tail rank by $M+1$ and counting
ties into the right tail gives $\Pr(\hat p_M\le\alpha\mid\mathcal
C)\le\alpha$. (ii) For fixed $S_0$ the indicators $\mathbf 1\{T(S_m)\ge
T(S_0)\}$, $m\ge1$, are i.i.d.\ Bernoulli$(q_\theta)$ given $\mathcal C$; the
strong law of large numbers yields the limit, which does not invoke $H_0^\theta$.
\end{proof}

Part~(i) is the finite-sample exactness of a Monte Carlo randomization test:
drawing $M$ schedules rather than enumerating $\mathbb{Q}_\theta$ preserves exact
level-$\alpha$ validity \cite{Dwass1957,Hope1968}, the modern group-theoretic
account is \cite{HemerikGoeman2018}, and the $+1$ in numerator and denominator is
what makes $\hat p_M$ valid at finite $M$---the naive $\sum\mathbf
1\{\cdot\}/M$ is downward biased and can be invalidly zero
\cite{PhipsonSmyth2010}. Because the null is generated by a constrained,
dependence-preserving resampling of one realized path rather than by i.i.d.\
draws, the relevant guarantee is that of \cite{BesagClifford1989}, which is exact
for precisely this conditional-on-the-data construction. The dependence of the
rejection probability on $M$ is honestly nonmonotone: on each plateau where
$k_\alpha=\lfloor\alpha(M+1)\rfloor-1$ is constant the rejection probability
declines, jumping upward only when $\alpha(M+1)$ crosses an integer, so the
finite-$M$ power traces a sawtooth around its $M\to\infty$ limit
\cite{Hope1968}; at $M=5{,}000$ and $\alpha=0.05$ this budget is far from binding.

\subsection{Consistency and the minimum detectable edge}

Validity is the floor; the test must also detect a genuine edge as evidence
accumulates. Evidence here accumulates in the trade count $N$, not in the length
of the price path, because $\mathbb{Q}_\theta$ rearranges a fixed set of $N$
trade contributions on a fixed path.

\begin{proposition}[Consistency and the $N^{-1/2}$ detection rate]
\label{prop:consistency}
Consider a per-trade location-shift alternative in which the realized placement
confers a common standardized entry edge $\zeta=\mu_1/\sigma_1>0$ on each of the
$N$ trade contributions to $\log(1+CR)$, relative to a $\mathbb{Q}_\theta$-mean of
zero. Let $Z_N=(T(S_0)-\mu_{\mathbb Q})/\sigma_{\mathbb Q}$ be the standardized
statistic, with $\mu_{\mathbb Q},\sigma_{\mathbb Q}$ the conditional
$\mathbb{Q}_\theta$ mean and standard deviation. Assume the conditional
finite-population/weak-dependence central limit regularity of
Remark~\ref{rem:clt}. Then $Z_N=\sqrt N\,\zeta+O_p(1)$, so the realized
$\mathbb{Q}_\theta$-percentile tends to one and the level-$\alpha$ test is
consistent as $N\to\infty$. Equivalently, the minimum edge detectable at fixed
power shrinks as $\zeta_{\min}\asymp N^{-1/2}$.
\end{proposition}

\begin{proof}[Proof sketch]
Write $\log(1+CR)=\sum_{j=1}^{N}\xi_j$, where under $\mathbb{Q}_\theta$ the $\xi_j$
are the $N$ trade log-contributions evaluated at a uniformly randomized placement
of the fixed duration/gap template on the fixed path. By
Remark~\ref{rem:clt}, $(T(S_0)-\mu_{\mathbb Q})/\sigma_{\mathbb Q}\to_d N(0,1)$
under $H_0^\theta$, and the location shift adds $N$ independent mean increments of
size $\mu_1$ with per-trade standard deviation $\sigma_1$, so the centered,
$\sigma_{\mathbb Q}$-scaled statistic is displaced by $N\mu_1/\sigma_{\mathbb
Q}=\sqrt N\,(\mu_1/\sigma_1)\,(\sqrt N\,\sigma_1/\sigma_{\mathbb Q})$, which is
$\sqrt N\,\zeta+O_p(1)$ once $\sigma_{\mathbb Q}\asymp\sqrt N\,\sigma_1$. The tail
probability $q_\theta$ then tends to zero, and Proposition~\ref{prop:q-validity}(ii)
makes $\hat p_M$ track it. Fixing the noncentrality $\sqrt N\,\zeta$ at the value
delivering a target power inverts to $\zeta_{\min}\asymp N^{-1/2}$.
\end{proof}

\begin{remark}[The correct limit theorem]
\label{rem:clt}
The Gaussian reference for $Z_N$ is \emph{not} an i.i.d.\ Lyapunov statement.
Under $\mathbb{Q}_\theta$ the summands $\xi_j$ are a gap-permutation of one fixed,
autocorrelated return path, hence neither independent nor identically
distributed; the relevant tool is a finite-population/sampling-without-replacement
or weak-dependence central limit theorem carrying a \emph{no-dominant-term}
(Lindeberg-type, $\max_j\mathrm{Var}(\xi_j)/\sum_j\mathrm{Var}(\xi_j)\to0$)
condition. This condition can fail: when one or a few trades dominate the
return budget, the normal approximation degrades and the standardized statistic
should not be read as pivotal. We therefore treat $\mathrm{RCSI}_z$ as a
descriptive, approximately pivotal effect size only; the test's validity rests on
the exact exchangeability of Proposition~\ref{prop:q-validity}(i), not on the
Gaussian tail. The adequacy of the approximation is then an empirical matter,
confirmed by the Kolmogorov--Smirnov uniformity tests of
Section~\ref{sec:simulation-validation} and the controlled-strength power
calibration of Section~\ref{sec:positive-control}.
\end{remark}

Two consequences match the empirical record. Low-frequency strategies are weakly
identified: Squeeze Breakout's $N=4$ admits only a large $\zeta_{\min}$, which is
why its near-threshold $p$-value is read as exploratory rather than as evidence of
absence. And because power is governed by $N$, a non-rejection at large $N$ is
informative while a non-rejection at small $N$ is not.

\subsection{Local power and relative efficiency}

The $N^{-1/2}$ boundary of Proposition~\ref{prop:consistency} is sharp: at exactly
that rate the test has nondegenerate, computable power.

\begin{proposition}[Local power against placement alternatives]
\label{prop:local-power}
Along the Pitman sequence $\zeta_N=h/\sqrt N$ for fixed $h\ge0$, under the
regularity of Remark~\ref{rem:clt}, $Z_N\to_d N(h,1)$, and the one-sided
level-$\alpha$ test has asymptotic power
\[
\beta(h)=\Phi\big(h-z_{1-\alpha}\big),
\]
with $z_{1-\alpha}$ the standard normal quantile. Hence $\beta(0)=\alpha$,
$\beta(z_{1-\alpha})=\tfrac12$, and at $\alpha=0.05$ power reaches $0.80$ at
$h=z_{0.95}+z_{0.80}\approx2.49$.
\end{proposition}

\begin{proof}
Under the contiguous sequence the standardized statistic converges to $N(h,1)$,
so $\Pr(Z_N\ge z_{1-\alpha})\to1-\Phi(z_{1-\alpha}-h)=\Phi(h-z_{1-\alpha})$. The
three values follow by substitution; $\Phi(z_{0.80})=0.80$ gives the stated
$h\approx2.49$.
\end{proof}

\begin{proposition}[Equal efficiency against placement, orthogonality to
exposure]
\label{prop:orthogonality}
Let SP denote the structure-preserving test and let RC/SPA denote the
single-strategy specialization of White's Reality Check \cite{white2000} and
Hansen's SPA test \cite{hansen2005}---studentized bootstrap tests of mean trade
return. Along the placement Pitman sequence of Proposition~\ref{prop:local-power},
SP and RC/SPA share the noncentrality $h$, so the asymptotic relative efficiency
of SP against placement alternatives is one. Along the orthogonal \emph{pure
structural-exposure} sequence---in which a strategy is profitable but its
placement carries no edge, so the realized statistic remains a typical draw under
$\mathbb{Q}_\theta$---SP has noncentrality \emph{zero}, whereas RC/SPA have
strictly positive noncentrality and reject.
\end{proposition}

\begin{proof}[Proof sketch]
Decompose realized profit into an exposure component, the return any structurally
identical schedule earns on the path, and a placement component, the excess of the
realized schedule over its $\mathbb{Q}_\theta$ peers. SP studentizes by the
conditional $\mathbb{Q}_\theta$ standard deviation, under which the exposure
component is held fixed across resamples and contributes nothing to the
conditional distribution: a profitable-but-untimed schedule is, by construction,
a central draw, so the noncentrality of SP is the placement component alone,
which equals $h$ along the placement sequence and $0$ along the exposure
sequence. RC/SPA studentize the realized mean trade return against a zero
benchmark; their noncentrality is the \emph{total} mean return and so loads on
both components. Along the exposure sequence the placement component vanishes but
the exposure component does not, giving RC/SPA strictly positive noncentrality
while SP carries zero. Equality along the placement sequence follows because there
the two noncentralities coincide.
\end{proof}

Orthogonality is the precise sense in which SP and the return-based tests answer
different questions: SP is, by construction, orthogonal to the exposure direction
that drives the competitors' rejections. The finite-sample shadow of this is the
controlled experiment of Section~\ref{sec:competitor}: in the structural-exposure
world, where a profitable strategy times its entries at random, SP holds its
nominal size at $0.050$ while Reality Check and SPA reject at $0.660$ and $0.638$
respectively---rejecting a strategy with no timing skill roughly two times in
three---and all four tests have full power ($1.000$) against genuine entry skill.
The methods do not differ in power against a real edge; they differ in what they
treat as an edge.

\subsection{The measure family and partial identification}

The preceding results condition on a single measure $\mathbb{Q}_\theta$. But the
boundary between ``structure'' (conditioned away) and ``placement'' (tested) is a
modeling choice, and different defensible choices define different nulls. We make
the conditioning measure itself part of the inferential object and ask what is
identified once it is acknowledged to be a choice.

\begin{definition}[Admissible measure class and identified verdict set]
\label{def:measure-family}
Let $\mathfrak Q=\{\mathbb{Q}_\theta:\theta\in\Theta\}$ be the family of
structure-preserving measures that all fix $(\mathcal S,\{P_t\})$ and differ only
in which additional features of the realized schedule they hold fixed---the
gap-permutation measure (conditioning on the geometry $\mathcal S$ alone, the
minimal element), the leading-gap restriction (forbidding placement in the
universal warm-up region), and the context-matched measure (restricting entries
to bars of comparable trailing volatility). Each $\mathbb{Q}_\theta$ defines a null
$H_0^\theta$ and an estimand $q_\theta=\mathbb{Q}_\theta\big(T(S)\ge
T(S_0)\big)$. Fix an admissible subclass $\Theta_0\subseteq\Theta$. The data
identify not a scalar ``skill'' but the \emph{verdict set}
$\{q_\theta:\theta\in\Theta_0\}$, equivalently the bounds
$[\,\min_{\theta\in\Theta_0}q_\theta,\ \max_{\theta\in\Theta_0}q_\theta\,]$.
\end{definition}

Point identification holds only \emph{given} $\theta$: $q_\theta$ is a
well-defined functional of the data, but ``the timing skill'' is not, because no
$\theta$ is privileged by the data. The estimand is therefore the verdict set,
and the natural reduction to a single claim is invariance across $\Theta_0$.

\begin{definition}[Measure-invariant null and skill claim]
\label{def:invariant}
The \emph{measure-invariant null} is the intersection
$H_0^{\cap}=\bigcap_{\theta\in\Theta_0}H_0^\theta$: there is some admissible
measure under which the realized placement is exchangeable with random placement.
\emph{Measure-invariant entry-placement skill} is its complement: rejection of
$H_0^\theta$ for \emph{every} $\theta\in\Theta_0$.
\end{definition}

The claim worth making is the complement of $H_0^{\cap}$. Testing it is an
intersection--union problem: the alternative (skill under every admissible
measure) is an intersection of one-sided alternatives, so a valid test rejects
only at the worst-case, least-favorable admissible measure.

\begin{proposition}[A valid test for measure-invariant skill]
\label{prop:iut}
Let $\hat p_M^\theta$ be the Monte Carlo $p$-value of
Proposition~\ref{prop:q-validity} computed under $\mathbb{Q}_\theta$, and define
\[
p^{\mathrm{inv}}=\max_{\theta\in\Theta_0}\hat p_M^\theta .
\]
Rejecting measure-invariant skill's null---i.e.\ rejecting $H_0^{\cap}$---when
$p^{\mathrm{inv}}\le\alpha$ is a level-$\alpha$ test: under $H_0^{\cap}$,
$\Pr(p^{\mathrm{inv}}\le\alpha\mid\mathcal C)\le\alpha$.
\end{proposition}

\begin{proof}
$H_0^{\cap}$ holds iff $H_0^{\theta^\ast}$ holds for at least one
$\theta^\ast\in\Theta_0$. On the event $H_0^{\theta^\ast}$,
Proposition~\ref{prop:q-validity}(i) gives
$\Pr(\hat p_M^{\theta^\ast}\le\alpha\mid\mathcal C)\le\alpha$. Since
$p^{\mathrm{inv}}=\max_\theta\hat p_M^\theta\ge\hat p_M^{\theta^\ast}$, the event
$\{p^{\mathrm{inv}}\le\alpha\}$ implies $\{\hat p_M^{\theta^\ast}\le\alpha\}$,
whence
$\Pr(p^{\mathrm{inv}}\le\alpha\mid\mathcal C)\le\Pr(\hat
p_M^{\theta^\ast}\le\alpha\mid\mathcal C)\le\alpha$.
This is the intersection--union test of \cite{LehmannRomano2005}: the maximum
$p$-value over an admissible class controls level at the least favorable element,
with no correction for $|\Theta_0|$.
\end{proof}

The test is honestly conservative: its size at the boundary of $H_0^{\cap}$ is
typically below $\alpha$ because only one $\theta^\ast$ need be null, and it makes
no claim of being most powerful. Its virtue is that ``significant under every
defensible measure'' is exactly the claim a referee can trust without adjudicating
which $\theta$ is correct. Two logical facts bound its interpretation, and we
state them precisely because they discipline what the design can and cannot
deliver.

\begin{proposition}[Necessity without sufficiency; non-necessity]
\label{prop:logic}
Let $\Theta_0$ be admissible. Then:
\emph{(i)} rejection of $H_0^\theta$ for every $\theta\in\Theta_0$ is
\emph{necessary} for a measure-invariant skill claim, but \emph{not sufficient}
for genuine entry-placement skill: a confound that displaces $T(S_0)$ in the same
direction under every $\mathbb{Q}_\theta\in\mathfrak Q$ survives the intersection
and is reported as ``skill'' though it is not placement;
\emph{(ii)} unanimous rejection is \emph{not necessary} for genuine
state-conditional skill: a measure $\mathbb{Q}_{\theta'}\in\Theta_0$ can absorb a
real state-conditional edge into its conditioning set, so that
$H_0^{\theta'}$ is not rejected even though placement genuinely carries
information.
\end{proposition}

\begin{proof}
(i) Necessity is immediate from Definition~\ref{def:invariant}: the claim is the
conjunction. For non-sufficiency, suppose a feature $C$ correlated with high
realized return is held fixed (hence not randomized) by every
$\mathbb{Q}_\theta\in\mathfrak Q$. Then $C$ contributes identically to $T(S_0)$ and
to each $T(S)$ under every measure; if $C$ also displaces $T(S_0)$ upward relative
to the resampled placements---because the realized schedule's interaction with the
common conditioned feature is favorable---then $q_\theta$ is small for all
$\theta$, so $H_0^\theta$ is rejected for all $\theta$, yet the displacement is
attributable to $C$, a confound shared by the whole class, not to placement. The
intersection cannot separate a placement effect from a common-to-all-measures
confound. (ii) Take a strategy whose edge is genuinely conditional on a market
state $V$ (e.g.\ trailing-volatility regime). The context-matched measure
$\mathbb{Q}_{\theta'}$ restricts resampled entries to bars matching the realized
$V$-context; under $\mathbb{Q}_{\theta'}$ the realized schedule competes only
against placements sharing its state, so the state-conditional edge is held fixed
rather than tested, and $T(S_0)$ is a central draw: $q_{\theta'}$ is not small and
$H_0^{\theta'}$ is not rejected, even though the strategy possesses genuine
state-conditional timing skill. Hence unanimity fails despite real skill.
\end{proof}

Part~(ii) is not hypothetical: Section~\ref{sec:schedule-sensitivity} exhibits
exactly such a case, in which four strategies reject under the context-matched
measure while the neutral gap-permutation measure rejects none, so a genuine
state-conditional volatility-timing edge cannot here be separated from a matched-pool
artifact. We therefore present measure invariance (Definition~\ref{def:invariant})
as the strongest claim the design supports, not as a characterization of skill:
it is the largest claim that survives the analyst's choice of the
structure/placement partition. We adopt the gap-permutation measure as the
minimal element of $\mathfrak Q$---it adds no contextual restriction beyond
$\mathcal S$---as the primary null, without asserting it is the unique admissible
choice, and report the verdict set across $\Theta_0$ rather than a single number.

% NOTE TO EDITOR: I have NOT asserted that p^{inv} is exact at the boundary
% of H_0^cap (it is conservative, as proved), nor that the verdict-set bounds
% are sharp in the econometric partial-identification sense — sharpness would
% require characterizing the full set of placement effects consistent with the
% data across ALL admissible theta, which the paper's three-measure family does
% not exhaust. Both are stated as conservative / set-valued rather than sharp,
% which is the defensible claim given the construction.
